\documentclass{article}
\usepackage[utf8]{inputenc}
\usepackage[margin=2cm]{geometry}
\usepackage{amsmath}
\usepackage{tikz}
\usepackage{pgfplots}
\pgfplotsset{compat=1.18}
\usepackage{needspace}
\usepackage{xcolor}
\usepackage{sectsty}
\usepackage{helvet}
\renewcommand{\familydefault}{\sfdefault}
\sectionfont{\normalfont\fontsize{16}{20}\selectfont}
\setlength{\parindent}{0pt}
\pagestyle{empty}
\begin{document}
\large

\section*{Teoria: Ecuaciones Lineales}

\vspace{3\baselineskip}
\begin{center}
\textbf{Forma general}

\vspace{0.8cm}
\fbox{\parbox{0.7\textwidth}{\centering \vspace{0.1cm} {\Large\color{black} \[ ax + b = 0 \quad \Rightarrow \quad x = \frac{-b}{a} \] } \vspace{0.1cm}}}
\end{center}
\vspace{0.6cm}
\begin{itemize}
\item {\Large \(a\)}: coeficiente del termino con \(x\)
\item {\Large \(x\)}: la incognita a despejar
\item {\Large \(b\)}: termino independiente
\end{itemize}
\vspace{2\baselineskip}

\needspace{4\baselineskip}
\hrule
\vspace{0.4cm}
Una ecuación lineal es aquella donde la incógnita aparece elevada solo a la
primera potencia, sin multiplicarse por sí misma ni aparecer dentro de raíces.
Su forma general es \(ax + b = 0\), con \(a \neq 0\).

Resolverla consiste en despejar x usando operaciones inversas: se resta b
de ambos lados, y luego se divide por a, manteniendo la igualdad en todo momento.
\vspace{0.4cm}
\hrule
\vspace{1cm}

\bigskip
\needspace{6\baselineskip}
\subsection*{Ejemplo 1: Caso directo}
\vspace{0.4cm}
Resolver \(3x + 6 = 0\).
\vspace{0.6cm}
{\large \[ 3x = -6 \quad \Rightarrow \quad x = -2 \] }
\vspace{0.6cm}
La solución es \(x = -2\).

\bigskip
\needspace{6\baselineskip}
\subsection*{Ejemplo 2: Incógnita en ambos lados}
\vspace{0.4cm}
Resolver \(5x - 3 = 2x + 9\).
\vspace{0.6cm}
{\large \[ 5x - 2x = 9 + 3 \quad \Rightarrow \quad 3x = 12 \quad \Rightarrow \quad x = 4 \] }
\vspace{0.6cm}
La solución es \(x = 4\).

\bigskip
\needspace{6\baselineskip}
\subsection*{Ejemplo 3: Con fracciones}
\vspace{0.4cm}
Resolver \(\frac{x}{2} + 3 = 7\).
\vspace{0.6cm}
{\large \[ \frac{x}{2} = 4 \quad \Rightarrow \quad x = 8 \] }
\vspace{0.6cm}
La solución es \(x = 8\).

\bigskip
\needspace{6\baselineskip}
\subsection*{Ejemplo 4: Sin solución (caso especial)}
\vspace{0.4cm}
Resolver \(2x + 5 = 2x + 8\).
\vspace{0.6cm}
{\large \[ 2x - 2x = 8 - 5 \quad \Rightarrow \quad 0 = 3 \] }
\vspace{0.6cm}
Esta igualdad es falsa, por lo tanto la ecuación no tiene solución.

\end{document}
